\documentclass[11pt,letterpaper]{article}
\usepackage[T1]{fontenc}
\usepackage[utf8]{inputenc}
\usepackage{lmodern}
\usepackage[margin=0.86in]{geometry}
\usepackage{microtype}
\usepackage{booktabs,longtable,tabularx,array}
\usepackage{enumitem}
\usepackage{xcolor}
\usepackage{hyperref}
\usepackage{xurl}
\definecolor{DocBlue}{HTML}{16435F}
\hypersetup{colorlinks=true,urlcolor=DocBlue,linkcolor=DocBlue}
\setlength{\parindent}{0pt}
\setlength{\parskip}{0.54em}
\setlist[itemize]{leftmargin=*,itemsep=3pt}
\begin{document}
\begin{center}
{\LARGE\bfseries Merge-Blocker Governance and Root-Cause Analysis}\par
\vspace{5pt}{\large Separate checks, rulesets, review obligations and pending dismissals}\par
\vspace{9pt}{\small Engineering reference | Proposed design | Not a deployment or approval record}\par
\end{center}
\tableofcontents
\clearpage
\section{Control taxonomy}
GitHub code scanning ruleset merge protection is distinct from required status checks. A dismissed code scanning alert does not automatically satisfy missing required reviews, check failures, merge-queue requirements or scanner completion.

\begin{itemize}
\item REQUIRED\_REVIEW: CODEOWNERS, AppSec teams, change requests or branch approval rules.
\item CODE\_SCANNING\_RULESET: tool, severity threshold, configuration and analysis state.
\item REQUIRED\_CHECK: job success, failure, missing name, skipped/cancelled or running.
\item DISMISSAL\_REQUEST: delegated request awaiting independent review.
\item MERGE\_QUEUE and CONFLICT: queue-specific semantics, stale heads and branch policies.
\end{itemize}

\section{Read-only diagnostic order}
\begin{itemize}
\item 1. Read PR base/head and mergeability; capture timestamp, ETag/commit.
\item 2. Inspect effective rulesets and branch-protection outcomes when permissions permit.
\item 3. Inspect check suites, required context names, and their exact states.
\item 4. Inspect reviews, CODEOWNERS, dismissal request state and approval authority.
\item 5. Inspect configured scanners, analyses and alert instances for head/diff relevance.
\item 6. Return distinct blockers with evidence and unknown-state diagnostics; never infer that alert counts alone block merge.
\end{itemize}

\section{Findings and limits}
Code scanning merge protection can block when required tool analysis is in progress or not configured. It does not inherently depend on a successful status check. Applicability to merge queues and Dependabot default-setup PRs must be verified against current platform limitations.

Not every repository equals an application. Monorepo components can have different owners, and CI relationships alone do not prove runtime exposure.


\section{Resolution routing}
\begin{itemize}
\item A real vulnerability requires remediation, compensating controls and authorized risk acceptance when permitted.
\item A suspected false positive needs source-to-sink and technical evidence, not merely a developer comment.
\item A missing or failed scan requires repairing pipeline/configuration, never dismissing a code finding.
\item A required AppSec PR approval requires the relevant reviewer or an authorized review-policy change.
\item A queue/conflict requires merge mechanics investigation rather than changing a security alert.
\end{itemize}

\section{Measurement contract}
Collect per-PR blocking category, first observed and resolved timestamps, repository criticality, review duration and denominator (eligible PRs). Avoid asserting causal improvements unless a representative pilot population and baseline are available.


\section*{Document status and provenance}
This document is a proposed reference design. Actual GitHub feature availability, organization policies, approvals, and runtime performance must be independently verified before adoption. The architecture inventory is not proof of vulnerability exploitability or correctness of a dismissal.
\end{document}
